\documentclass[11pt]{article}
\usepackage[margin=0.8in]{geometry}
\usepackage[]{xcolor}
\parindent=0pt

\title{CSE 4300: Exam 1}
\author{gordonbchen}
\date{}

\begin{document}
\maketitle


\section{P1}
32-bit addr space, 8kb page size, 32-bit entry, copy 1 word = 100ns. each proc runs for 100ms.

Find \% of CPU time to load page table.\\

13-bit offset, so 19-bit page table idx, $2^{19}$ page table entries.

$2^{19} * 1$ words = $2^{19}$ words, so $100ns * 2^{19} = 10^-7 * 10^6 / 2 = 10^-1 / 2 = 50 ms$\\

So 50\% CPU time.


\section{P2}
64kb virtual address space. 16kb physical address space. 4kb page size.\\

64kb VA space, so 16 bit virtual address. 16kb RAM, so 14 bit virtual address.\\

4kb page size, so 12 bit offset, 4 bit page table idx.\\

4 bit virtual page idx $\to$ 2 bit physical page idx.


\section{P3}
64kb program size. 2kb page size. Time to load program HDD $\to$ RAM?

Magnetic HDD: 5ms seek time (get to correct track), 5ms rotation time (get right bit on track)

on avg have to wait half rotation to read/write.

1mb track.\\

5ms to read 1 mb. 5ms * 64/1024 to read 64kb.

then 5ms seek + 5/2 ms rotate + 5ms * 64/1024.


\section{P4}
amount of memory proportional to time between page faults

1micros for instruction w/out page fault, 2001micros for instruction w/ page fault (2ms for page fault)

60s to run program, 15000 page faults.

15000 * 2ms = 30 000 ms = 30s time on page faults, so 30s instructions.

if double memory, how much time?

half time on page fault, so 15s page faults, total 45 sec.


\section{P5}
Shortest process next: predict execution time.

Times; 20, 30, 40.

$a T_0 + (1-a) T_1, a = 1/2$

20, (20 + 30) / 2 = 25, (25 + 40) / 2 = 32.5
\end{document}
